\section{Discussion}

\subsection{Where the benefit of spatially aware wiring lies}

Taken together, the two unsupervised findings identify where the benefit of spatially
aware wiring actually lies, and they temper what can be expected of unsupervised
plasticity acting on it. Because the oriented fields already supply much of the structure
these tasks require, an unsupervised rule has relatively little left to contribute and a
good deal that it can disturb, so any rule that modifies these weights must recover more
structure than it removes before it improves on leaving them untouched. None of the
unsupervised rules we tested managed that, and the shortfall was consistent enough across
datasets and conditions that we understand it as a consequence of how much the
initialization already provides rather than a deficiency particular to any one rule.

That the erosion of the fields does not predict accuracy suggests orientation coherence
is better read as a description of what plasticity does to the representation than as an
account of why one configuration outperforms another. Two conditions with identical
coherence differ by three percentage points, and the supervised model reaches its best
results while retaining roughly seventy percent of the initial value, so we do not treat
retention as a quantity to be maximized.

\subsection{Supervision and the structural prior}

Supervision does not free the network from its structural assumption; it sharpens it.
Read by the same external probe in both phases, the advantage of oriented fields over
random connectivity of matched density grows from $+0.64$ to $+4.92$ percentage points on
MNIST and from $+3.48$ to $+5.31$ on KMNIST once the rule carries a teaching signal, while
on notMNIST the identical shift turns a $+3.24$ advantage into a $-2.64$ cost. Through its
own readout the supervised network reaches $94.6\%$ on MNIST and $75.5\%$ on KMNIST, the
two datasets on which the prior and the rule agree. We read this as the reward signal
amplifying whatever the initialization already encodes rather than replacing it: the prior
functions as an inductive bias in the ordinary sense, narrowing what the network can
represent in exchange for aligning it with a hypothesis about the input, so that the
supervised objective sharpens both the benefit and the cost of that trade.

What changes between the two phases is not the amount of plasticity but whether it is
directed. Holding the architecture and the prior fixed and setting the reward rate to
zero, R-STDP improves on frozen weights on all five datasets, by $8.9$ percentage points
on MNIST and $22.7$ on KMNIST, where every unsupervised rule we tested fell below the
same reference. We therefore read the phase-1 result as a statement about undirected
local plasticity rather than about plasticity as such: a rule that receives no
information about the task can only rearrange a prior that was already close to what the
task needed, while a rule that receives such information can improve on it.

We do not, however, claim that the supervised model is more accurate than the
unsupervised one. Six factors change between the two phases, so no single difference is
attributable to supervision, and the frozen comparison above is valid only within
phase~2. What supervision demonstrably adds is a decoder internal to the network, which
the unsupervised model lacks entirely.

\subsection{Scope of the abstention result}

Abstention is useful where the model is already accurate and its confidence is
calibrated, and it offers little where either condition fails. The separation between
correct and incorrect predictions weakens in step with accuracy across the five datasets,
so at a fixed accuracy target the usable fraction of the test set falls from most of
MNIST to a few percent of Fashion-MNIST and none of SVHN. Reporting abstention as a
general capability of the architecture would overstate what the mechanism provides; on
the present evidence it is a property of the MNIST-trained model rather than of the
architecture.

\subsection{Where novelty is represented}

The out-of-distribution result separates two things the abstention analysis treats as
one. Predictive entropy is a statement about the readout's output distribution, and it
identifies novel input poorly; the same trial, read as a point in the excitatory rate
space, separates novel from familiar input considerably better. The information was in
the representation throughout and the readout discarded it, which bounds the preceding
section rather than contradicting it: a calibrated confidence signal reports when the
network is likely to be wrong about a class, not when the input belongs to no class at
all.

The second observation is structural. We adopted the class-grouped layout so that reward
could be assigned to a fixed population per class, and it turns out to supply a novelty
signal at no additional cost, since familiar input concentrates activity in one group
while novel input spreads it across the ten and no parameters need to be fitted to
measure that. A choice made for credit assignment paying a second dividend elsewhere is
the same pattern we report for the receptive-field prior. What the network does not
provide is a model of the input distribution, and a linear subspace fitted to raw pixels
detects novelty better than any score we read from the network---which, on a benchmark of
centred grayscale images, is what we would expect.

\subsection{Limitations}

The prior is manipulated only under trace-STDP in the unsupervised phase, so we cannot
say whether receptive fields would help under triplet-STDP or under frozen weights. A
frozen random-weight condition would separate the contribution of the oriented geometry
from that of freezing the weights, and its absence is the main gap in the present design.
The supervised contrast also depends on which decoder reads the representation: on
Fashion-MNIST the learned readout and the external probe disagree in significance, and
the two differ in magnitude on every dataset, so each statement of that contrast has to
name its decoder. The out-of-distribution analysis rests on one trained network at a
single seed, and the two feature-space scores it compares were selected after seeing the
probes, so we report it as a characterisation rather than a benchmarked result.

\subsection{Conclusion and outlook}

Biological realism is not a single property that a network either has or lacks, and our
results separate its components rather than endorsing or rejecting it as a whole.
Spatially structured connectivity pays where the data carry the structure it assumes and
costs accuracy where they do not; undirected local plasticity improves on that structure
on none of the five datasets; and the same plasticity, once modulated by reward, improves
on it on all of them while amplifying the prior's advantage and its disadvantage alike.
The resulting model is accurate enough on MNIST for its own confidence to support a
usable abstention policy and not accurate enough elsewhere for that policy to mean much,
which we take as the honest scope of the approach at this scale.

Two directions follow from this. The sign of the receptive-field effect should be
predictable from a statistic of the images rather than recovered after training, which
would turn the choice of prior into a design decision. The second concerns the rule
itself: preservation of the prior and task performance came apart in our comparisons,
since the rule that preserved the most was not the one that performed best, and the rule
that performed best was the one carrying a teaching signal. Combining the two properties
is therefore more promising than trading one against the other.
